\documentclass[11pt]{article}
\usepackage{fontspec}
\usepackage[french,english]{babel}
\usepackage{amsmath}
\usepackage{amssymb}
\usepackage{amsthm}
\usepackage{mathrsfs}
\usepackage{tikz-cd}
\usepackage{enumitem}

\newcommand{\printedpage}[1]{\marginpar{\footnotesize p.~#1}}
\newcommand{\msp}[1]{\marginpar{\footnotesize [#1]}}
\newcommand{\upi}{\underline{\pi}}
\newcommand{\uCh}{\underline{\mathrm{Ch}}}
\newcommand{\Ob}{\operatorname{Ob}}
\newcommand{\Hom}{\operatorname{Hom}}
\newcommand{\bDelta}{\boldsymbol{\Delta}}

\begin{document}

\setcounter{footnote}{65}

\printedpage{68}
% head: A. Grothendieck : Dérivateurs, Ch. XII, Édition des Universités de Montpellier II et Paris VII

\noindent\textbf{Proposition 9.}

\begin{enumerate}[label=\textit{\alph*)}]
\item \textit{Every object $V$ of} Cat \textit{having a final object is totally $W$-aspheric - hence every presheaf of sets $H$} representable \textit{on a category $U$ (representability meaning precisely that $U/H$ has a final object) is totally $W$-aspheric.}

\item (\footnote{!! Probably \textit{false} as it stands, cf. correction on the next page.}) \textit{(Suppose} W(7 bis) \textit{is satisfied.) Let $V$ be in} Cat, \textit{$V'$, $V''$ two opens with union $V$. If $V'$, $V''$, $V'\cap V''$ are totally $W$-aspheric, the same holds for $V$.}
\end{enumerate}

\noindent\textsc{Proof} of a). Take $V$ with a final object and $f : V \longrightarrow V'$, whence the factorization
\[
\begin{tikzcd}[column sep=small, row sep=small]
 & \tilde{V} = \upi_0(V/V') \arrow[dr, "\pi"'] & \\
V \arrow[ur, "\tilde{f}"] \arrow[rr, "f"'] & & V'
\end{tikzcd}
\]
of $f$. One proves without difficulty (\footnote{Here is a simple proof: For every $V$ in Cat we have $e_V = \varinjlim_{v \text{ in } V} v$ (direct limit of the inclusion functor of $V$ into $V^\wedge$), for $\forall\, x$ in $V$
\[
\varinjlim_{v \text{ in } V} \Hom(x,v) \simeq \pi_0(x\backslash V) = e.
\]
Since $f_!$ commutes with $\varinjlim$ , we therefore have $f_!(e_V) = \varinjlim_{v \text{ in } V} f(v)$. If $v_0$ is a final object of $V$, this shows that $f_!(e_V) = f(v_0)$, q.e.d.}) (a matter of cofinality) the following:

\noindent\textbf{Corollary.} \textit{The image under $\tilde{f}$ of a final object of $V$ is a final object of $\tilde{V} = \upi_0(V/V')$.}

Since an object of Cat with a final object is $W$-aspheric by W(3), this establishes a).

Let us prove b). I feel like saying that if we have $f : V \longrightarrow W$ in Cat, then $\upi_0(V/W)$ is a union of two opens isomorphic to $\upi_0(V'/W)$, $\upi_0(V''/W)$, whose intersection is isomorphic to $\upi_0(V'\cap V''/W)$. This would mean neither more nor
\msp{page 98}
less than this: $f_!$ transforms the monomorphisms $V' \longrightarrow V$, $V'' \longrightarrow V$ (interpreted as arrows in $V^\wedge$, with target $e_V$) into monomorphisms, and the fiber product $V'\cap V'' = V' \times_V V''$ into the corresponding fiber product in $W^\wedge$. But the functor $f_!$ commutes with arbitrary $\varinjlim$ (being a left adjoint), and I see no a priori reason why it should commute (say) with fiber products (which would also imply that it transforms monomorphisms into likewise). If I remember correctly, this is ultrafalse in general. So I now very much doubt that statement b) is true, without an additional hypothesis. To ensure monomorphy, one thinks of requiring the existence of \textit{retractions} of $V$ onto $V'$, $V''$ and of $V'$ or $V''$ onto $V'\cap V''$.

I have just checked this, and give the
% foot: Version 7 décembre 2002
\printedpage{69}
% head: A. Grothendieck : Dérivateurs, Ch. XII, Édition des Universités de Montpellier II et Paris VII
\textbf{Corollary 1} (correction of \texttt{[the]} prop. 9, part b) (\footnote{NB There is a slightly more general statement, starting from a category $U$ and a diagram
$\begin{tikzcd}[ampersand replacement=\&, row sep=small]
V'\cap V'' \arrow[r] \arrow[d] \& V' \arrow[d] \\
V'' \arrow[r] \& V
\end{tikzcd}$
\textit{in} $U^\wedge$: if $V'$, $V''$, $V'\cap V''$ are totally $W$-aspheric, the same holds for $V$.}): \textit{The statement is valid if one assumes that there exists a} retraction \textit{$r'_0$ of $V''$ onto $V'\cap V''$, and similarly a retraction $r''_0$ of $V'$ onto $V'\cap V''$.} (\footnote{Corrected corollary also false - the retractions are not arrows of $V^\wedge$, and do not pass into $p_!$.}) % [resolved: see DER_XIIb-a_corrections.csv]
\[
\begin{tikzcd}
V'\cap V'' \arrow[r, "i''_0"'] \arrow[d, "i'_0"] & V' \arrow[d, "i'"] \arrow[l, dotted, bend right=50, "r''_0"'] \\
V'' \arrow[r, "i''"] \arrow[u, dotted, bend left=50, "r'_0"] & V \arrow[u, dotted, bend right=50, "r'"'] \arrow[l, dotted, bend left=50, "r''"]
\end{tikzcd}
\]

(\textbf{NB} We will see later that the existence of \textit{one} of these retractions suffices to ensure the conclusion.)

\textsc{Proof}. Since $V$ is the amalgamated sum of $V'$ and $V''$, there exists a retraction $r'$ of $V$ onto $V'$ (given by $\mathrm{id}_{V'}$ on $V'$, by $i''_0r'_0$ on $V''$) which induces $r'_0$ on $V''$, and similarly there exists
\msp{page 99}
a retraction $r''$ of $V$ onto $V''$, inducing $r''_0$ on $V'$. The existence of these retractions already ensures that for every functor $f : V \longrightarrow W$, the transform of the diagram
\[
\begin{tikzcd}[row sep=large]
V'\cap V'' \arrow[r] \arrow[d] & V' \arrow[d] \\
V'' \arrow[r] & V
\end{tikzcd}
\]
(interpreted as a diagram in $V^\wedge$) consists of monomorphisms (\footnote{\textbf{false,} already in the case where $V = \Phi$, $V'$, $V''$ its two canonical opens isomorphic to $\bDelta^1$, and $V'{\cap}V'' = e$.}) % [sic] intérprété -> interprété, misspelt accent
\[
\begin{tikzcd}[row sep=large]
f_!(V'\cap V'') \arrow[r, hook] \arrow[d, hook] & f_!(V') \arrow[d, hook] \\
f_!(V'') \arrow[r, hook] & f_!(V),
\end{tikzcd}
\]
and moreover that $f_!(V') \cap f_!(V'')$ (which contains $f_!(V'\cap V'')$) is in fact \textit{equal to it.} It follows that we have a diagram both cartesian and cocartesian of open immersions in Cat
% foot: Version 7 décembre 2002
\printedpage{70}
% head: A. Grothendieck : Dérivateurs, Ch. XII, Édition des Universités de Montpellier II et Paris VII
\[
\begin{tikzcd}[row sep=large]
W/f_!(V'\cap V'') \arrow[r, hook] \arrow[d, hook] & W/f_!(V') \arrow[d, hook] \\
W/f_!(V'') \arrow[r, hook] & W/\Bigl(f_!(V) \overset{\text{def}}{=} f_!(e_V)\Bigr),
\end{tikzcd}
\]
which is also identified with
\[
\begin{tikzcd}[row sep=large]
\upi_0(V'\cap V''/W) \arrow[r, hook] \arrow[d, hook] & \upi_0(V'/W) \arrow[d, hook] \\
\upi_0(V''/W) \arrow[r, hook] & \underline{\pi_0}(V/W).
\end{tikzcd}
\]

\msp{page 100}
The hypothesis on $V'$, $V''$, $V'\cap V''$ implies that the three vertices of the square other than $\upi_0(V/W)$ are $W$-aspheric, hence by W(7 bis) the same holds for the latter, q.e.d.

\textbf{Corollary 2} (\footnote{More generally, if $I$ is a category \textit{ordered and finite,} having an initial object, then}). \textit{Assuming} W(7 bis)\textit{, the category} $\Psi = \begin{tikzcd}[ampersand replacement=\&, row sep=small, column sep=small, baseline=(B.base)]
 \& a \\
|[alias=B]| 0 \arrow[ur] \arrow[dr] \& \\
 \& b
\end{tikzcd}$ \textit{(or $0 \longleftarrow 1 \longrightarrow 2$) is totally $W$-aspheric} (\footnote{\textbf{False,} $\Psi$ is \textit{not} totally $W$-aspheric.}).

One applies corollary 1 to $\Psi$ and to its two open subcategories, $\Psi/a$, $\Psi/b$ (isomorphic to $\bDelta^1$), and to their intersection $\{0\} \simeq \bDelta^0 = e$. Since these three categories have a final object (unlike $\Psi$ itself, which has only an initial object \dots), they are totally $W$-aspheric. The corollary applies, the existence of the retractions $r'_0$, $r''_0$ being obvious here.

It is especially part b) (corrected) of page 9 which ensures the existence of totally $W$-aspheric categories other than those having a final object (\footnote{I now conjecture that these are the only ones! (\footnotemark)}\footnotetext{False, there are the filtering categories (if $W$ satisfies the axiom of limits).}). It is true that conversely one should still make sure that a $W$-aspheric object $X$ of Cat is not \textit{always} totally $W$-aspheric! The counterexample is not so obvious after all -
\msp{page 101}
I spent hours before finding one and getting a little clarity! It suffices, in short, to exhibit an example of an arrow in Cat
\[
f : X \longrightarrow Y
\]
such that
% foot: Version 7 décembre 2002
\printedpage{71}
% head: A. Grothendieck : Dérivateurs, Ch. XII, Édition des Universités de Montpellier II et Paris VII
\begin{enumerate}[label=\alph*)]
\item $\underline{\pi_0}(f) \overset{\text{def}}{=} \upi_0(X/Y) \xrightarrow{\ \sim\ } Y$, i.e. the $\pi_0(y\backslash X)$ \texttt{[rather the} $y\backslash X$\texttt{]} are 0-connected (for $y \in \Ob Y$),
\end{enumerate}

and such that moreover

\begin{enumerate}[label=\alph*), start=2]
\item $X$ is $W$-aspheric, and $Y$ not.
\end{enumerate}

% a vertical bar is printed in the left margin beside the next paragraph (lines "La difficulté ..." to "... le foncteur")
The difficulty is the following, when one looks for ``stupid'' counterexamples: it is that hypothesis a) already implies, when $X$ is $W_\infty$-aspheric (or only 1-connected) that $Y$ is 1-connected: the functor
\[
Y' \longmapsto X' = X \times_Y Y'
\]
from the category of étale coverings $Y'$ of $Y$ to that of étale coverings $X'$ of $X$ is fully faithful (immediate verification), whence the stated result. So a counterexample must call upon the existence of a nonzero $\pi_i(Y,y)$, with $i \geq 2$. This implies
\msp{page 102}
that there is no counterexample with categories of the form $Y = Z/F$ (deduced from a functor $q : X \longrightarrow Z$ by taking $F = q_!(e_X) = \pi_0(X/Z)$), where $Z$ is a category such as $\bDelta^1$ (?), $\Phi$, $\Psi$, or $\begin{tikzcd}[ampersand replacement=\&, sep=small, baseline=(M.base)]
 \& {} \& \\
{} \arrow[ur] \arrow[dr] \& |[alias=M]| \phantom{x} \& {} \arrow[ul] \arrow[dl] \\
 \& {} \&
\end{tikzcd}$ (having the homotopy type of the circle) etc., for the connected components of such categories $Y$ are $K(\pi,1)$'s - the higher $\pi_i$ ($i \geq 2$) are zero.

This difficulty seen, the counterexample becomes obvious. I start from a 1-connected object \textit{whatsoever} $Y$ of $X$ \texttt{[rather of} Cat\texttt{]}, and I take $y_0 \in \Ob Y$,
\[
X = \uCh(Y; y_0, -) \xrightarrow{\ f\ } Y.
\]
We know that $X$ is $W_\infty$-aspheric, and that $f$ is $W_\infty$-perfect, hence the inclusions
\[
X_y \longrightarrow y\backslash X
\]
are $W_\infty$-equivalences. Hypothesis a) that the $y\backslash X$ be 0-connected is therefore equivalent to the same hypothesis on the fibers $X_y$, or
\msp{page 103}
again on a single fiber ($f$ being perfect, hence a $W_\infty$-fibration, and $Y$ 0-connected), namely
\[
X_{y_0} \simeq \Omega(Y, y_0).
\]
Its 0-connectedness is equivalent to the hypothesis of 1-connectedness of $Y$ (so far we have used only that of 0-connectedness). So for a counterexample, it suffices to take $Y$ 1-connected and not $\infty$-connected. The simplest case is the one where $Y$ has the homotopy type of the sphere $S^2$, realized e.g. by the ordered category shown opposite \texttt{[below]} with 6 vertices.
% foot: Version 7 décembre 2002

\printedpage{72}
% head: A. Grothendieck : Dérivateurs, Ch. XII, Édition des Universités de Montpellier II et Paris VII
\[
\begin{tikzcd}[row sep=small]
e_2^+ \bullet & \bullet\, e_2^- \\
e_1^+ \bullet \arrow[ur] & \bullet\, e_1^- \arrow[ul] \\
e_0^+ \bullet \arrow[ur] & \bullet\, e_0^- \arrow[ul]
\end{tikzcd}
\]
\texttt{[moreover,} $e_0^+ \leq e_1^+ \leq e_2^+$\texttt{,} $e_0^- \leq e_1^- \leq e_2^-$\texttt{]} (\footnote{It is in short this example (up to topological homotopy) that I realize (without at first noticing it in the Appendix p. 10-12, but with finite categories).}). I summarize the result obtained.

\textbf{Proposition 10.} \textit{Let $Y$ be a small category which is} 1\textit{-connected ($y_0 \in \Ob Y$). Then} $\uCh(Y; y_0, -)$ \textit{is $W_\infty$-aspheric, but is not totally $W_\infty$-aspheric if $Y$ itself is not $\infty$-connected, i.e. $W_\infty$-aspheric.} (\footnote{But a more striking example is $\Psi = \begin{tikzcd}[ampersand replacement=\&, row sep=small, column sep=small, baseline=(A.base)]
 \& b \\
|[alias=A]| a \arrow[ur] \arrow[dr] \& \\
 \& c
\end{tikzcd}$ , which is $W_\infty$-aspheric, but not totally!}).

One may ask whether $Y$ $W_\infty$-aspheric implies that $\uCh(Y; y_0, -)$ is totally $W_\infty$-aspheric. I very much doubt it!

\msp{page 104}
Let us now return to the situation at the beginning of the section, diagram $(*)$, p.~87,
\[
\begin{tikzcd}
U \arrow[r, hook, "i"] \arrow[d, "p"'] & X \arrow[d, "q"'] & T = \{t\} \arrow[l, hook', dotted, "j"'] \arrow[d, dotted, "\mathrm{iso}"] \\
U' \arrow[r, hook, "i'"] & X' & T' = \{t'\}. \arrow[l, hook', dotted, "j'"']
\end{tikzcd}
\tag*{$(*)$}
\]
We can now give an example where $i$ is in $W$ (and even $W$-aspheric) without $i'$ being so:
\[
i \in W, \qquad i' \notin W. \tag*{$(**)$}
\]
One can even assume that $t$ is a final object of $X$ (which means that the object $H$ of $U^\wedge$ is isomorphic to $e_U$, the final object), and $t'$ a final object of $X'$ (i.e. $p_!(e_U) = e_{U'}$, i.e. the $u'\backslash U$ are 0-aspheric). In this case, since $X$, $X'$ are $W$-aspheric, the relations $(**)$ are equivalent to
\[
U \text{ is } W\text{-aspheric, and } U' \text{ not,} \tag*{$(***)$}
\]
which is precisely the situation we have just constructed (p. 101-103).

One would also like to find an example where in $(*)$ one has $p \in W$, $q \notin W$.
% foot: Version 7 décembre 2002

\end{document}
